\chapter{हिल्बर्ट समष्टियाँ}\label{ch:b3:hilbert}

\omterm{def:b3:hilbert:inner}{हिल्बर्ट समष्टि} वह बानाख समष्टि है जिसका मानक किसी \omterm{def:b3:hilbert:inner}{अंतःगुणन} से
आता है --- और यही एक अतिरिक्त संरचना अनंत विमा में यूक्लिडीय
ज्यामिति का लगभग सब कुछ लौटा देती है: लांबिक प्रक्षेप विद्यमान
होते हैं, प्रत्येक \omterm{def:b3:topology:continuity}{संतत} फलनिक किसी स्थिर सदिश के साथ \omterm{def:b3:hilbert:inner}{अंतःगुणन} ही
होता है (रीस), और प्रसामान्य लांबिक आधार प्रत्येक सदिश को
पाइथागोरसीय लेखा-जोखा वाली एक अभिसारी श्रेणी में फैला देते हैं
(\omterm{thm:b3:hilbert:parseval}{पार्सेवाल})। इस अध्याय का शिखर एक चुकाया गया ऋण है: त्रिकोणमितीय
निकाय $L^2$ का प्रसामान्य लांबिक आधार है, अतः \omterm{thm:b3:hilbert:parseval}{पार्सेवाल} की सर्वसमिका
\emph{प्रत्येक} वर्ग-समाकलनीय फलन के लिए लागू होती है --- वही कथन
जिसे दूसरा वर्ष केवल खंडशः $\mathcal C^1$ फलनों के लिए सिद्ध कर सका था। हम
लाक्स--मिलग्राम पर समाप्त करते हैं, जो अवकल समीकरणों के विचरण
उपागम की श्रमिक प्रमेयिका है।

सर्वत्र $H$ $K = \R$ या $\C$ पर एक सदिश समष्टि है।

\section{अंतःगुणन; प्रक्षेप प्रमेय}

\begin{definition}\label{def:b3:hilbert:inner}
\emph{अंतःगुणन}\index{अंतःगुणन} ऐसा प्रतिचित्रण $\langle
\cdot,\cdot\rangle \colon H\times H \to K$ है जो दूसरे
चर में रैखिक हो, जिसमें $\langle y, x\rangle =
\overline{\langle x, y\rangle}$ हो, और $x \neq 0$ के लिए $\langle x, x\rangle > 0$ हो। वह मानक
$\norm x = \langle x,
x\rangle^{1/2}$, \emph{कोशी--श्वार्ज़ असमिका} $\abs{\langle x, y\rangle} \leq \norm x\norm y$ (दूसरे वर्ष की उपपत्ति ---
विविक्तकर --- अपरिवर्तित रहती है), और \emph{समांतर चतुर्भुज नियम}
\[
\norm{x + y}^2 + \norm{x - y}^2 = 2\norm x^2 + 2\norm y^2 .
\]
प्रेरित करता है। \emph{हिल्बर्ट समष्टि}\index{हिल्बर्ट समष्टि} वह
अंतःगुणन समष्टि है जो इस मानक के लिए \omterm{def:b3:complete:complete}{पूर्ण} हो। उदाहरण: $\ell^2$ (\cref{pb:b3:banach:1})
और, सबसे मूलभूत, $\langle f, g\rangle = \int\bar fg\,\dd\mu$ के साथ $L^2(\mu)$ --- जो रीस--फिशर (\cref{thm:b3:lp:complete}) से \omterm{def:b3:complete:complete}{पूर्ण}
है; और अंतःगुणन कोशी--श्वार्ज़ ($p = q = 2$ पर $=$ होल्डर) से परिमित है।
\end{definition}

\begin{theorem}[किसी संवृत उत्तल समुच्चय पर प्रक्षेप]
\label{thm:b3:hilbert:projection}
मान लीजिए $C \neq \varnothing$ \omterm{def:b3:hilbert:inner}{हिल्बर्ट समष्टि} $H$ का संवृत \emph{उत्तल}
उपसमुच्चय है और $x \in H$। तब एक अद्वितीय $p_C(x)
\in C$ है जिसके
लिए\index{प्रक्षेप प्रमेय}
\[
\norm{x - p_C(x)} = d(x, C),
\]
और जो इस गुण से अभिलक्षित होता है: सभी $c \in C$ के लिए $\operatorname{Re}\langle x - p_C(x),\ c -
p_C(x)\rangle \leq 0$।
प्रतिचित्रण $p_C$ $1$-लिप्शिट्ज़ है।
\end{theorem}

\begin{proof}
मान लीजिए $d = d(x, C)$ और $\norm{x - c_n}
\to d$ वाले $(c_n) \subseteq C$। $x - c_n$ और $x - c_m$ पर समांतर चतुर्भुज:
\[
\norm{c_n - c_m}^2 = 2\norm{x - c_n}^2 + 2\norm{x - c_m}^2 -
4\,\bigl\|x - \tfrac{c_n + c_m}2\bigr\|^2
\leq 2\norm{x{-}c_n}^2 + 2\norm{x{-}c_m}^2 - 4d^2
\]
(उत्तलता मध्य बिंदु को $C$ में रख देती है): दायाँ पक्ष $0$ की ओर
जाता है, अतः $(c_n)$ कोशी है, और उसकी सीमा $p \in C$ (संवृत) $d$ प्राप्त
कर लेती है। अद्वितीयता: दो न्यूनीकारकों से उसी सर्वसमिका द्वारा
$\norm{p - p'}^2 \leq 2d^2 + 2d^2 - 4d^2 = 0$।

अभिलक्षण: $c \in C$, $t \in \intoc01$ के लिए सदिश $p + t(c - p) \in C$, अतः
\[
d^2 \leq \norm{x - p - t(c-p)}^2
= d^2 - 2t\operatorname{Re}\langle x - p, c - p\rangle +
t^2\norm{c-p}^2 ;
\]
अब $t \to 0^+$ से भाग दीजिए: $\operatorname{Re}\langle x - p, c -
p\rangle \leq 0$। विलोमतः यही असमिका $\norm{x -
c}^2 = \norm{x - p}^2 - 2\operatorname{Re}\langle x - p, c -
p\rangle + \norm{p - c}^2 \geq \norm{x-p}^2$ दे देती है।
लिप्शिट्ज़: प्रक्षेप $p, q$ वाले $x, y$ के लिए दोनों विचरण असमिकाएँ
(क्रमशः $c = q$ और $c = p$ के साथ) जोड़िए: $\operatorname{Re}\langle x - y - (p - q), p - q\rangle \geq
0$, अतः $\norm{p - q}^2 \leq \operatorname{Re}\langle x - y, p -
q\rangle \leq \norm{x - y}\norm{p - q}$।
\end{proof}

\begin{theorem}[लांबिक विघटन]
\label{thm:b3:hilbert:decomposition}
मान लीजिए $F$ $H$ की कोई \emph{संवृत उपसमष्टि} है। तब $p_F$
रैखिक है, सभी $x$ के लिए $x - p_F(x) \perp F$, और\index{लांबिक पूरक}
\[
H = F \oplus F^\perp,
\qquad F^\perp = \{y : \langle y, f\rangle = 0\ \forall f\in
F\},
\qquad (F^\perp)^\perp = F .
\]
किसी व्यापक उपसमष्टि के लिए $(F^\perp)^\perp = \bar F$; विशेष रूप से $F$ सघन है तभी
और केवल तभी जब $F^\perp = \{0\}$।
\end{theorem}

\begin{proof}
किसी उपसमष्टि के लिए $c =
p_F(x) \pm f$ ($f \in F$, दोनों चिह्न, और सम्मिश्र स्थिति
में $\iu f$) के साथ विचरण अभिलक्षण $\langle x - p_F(x), f\rangle = 0$ अनिवार्य कर देता है: अर्थात्
अवशेष $F$ के लांबिक है। विघटन $x = p_F(x) + (x
- p_F(x))$ जिसमें $F \cap F^\perp = \{0\}$ ($\langle y, y
\rangle = 0$); और $p_F$
की रैखिकता ऐसे विघटनों की अद्वितीयता से निकलती है (दोनों पक्ष उनमें
रैखिक हैं)। $(F^\perp)
^\perp \supseteq F$ सदैव; विलोमतः यदि $x \perp F^\perp$ हो, तो $x = f + g$ लिखिए: $g = x - f \in F^\perp$ और
$\langle g,
g\rangle = \langle x, g\rangle - \langle f, g\rangle = 0$: अतः $x
= f \in F$। किसी व्यापक उपसमष्टि $F$ के लिए: $F^\perp = \bar
F^{\,\perp}$ (\omterm{def:b3:hilbert:inner}{अंतःगुणन}
का \omterm{def:b3:topology:continuity}{सांतत्य}), अतः संवृत स्थिति से $(F^\perp)^\perp = \bar F$; और सघनता तभी और केवल तभी
जब $\bar F = H$, अर्थात् तभी और केवल तभी जब $F^\perp = 0$।
\end{proof}

\begin{theorem}[रीस निरूपण]\label{thm:b3:hilbert:riesz}
प्रत्येक \omterm{def:b3:topology:continuity}{संतत} रैखिक फलनिक $\varphi \in H'$ के लिए एक अद्वितीय $a \in H$ है जिसके
लिए\index{रीस निरूपण प्रमेय}
\[
\varphi(x) = \langle a, x\rangle \quad (x \in H),
\qquad \norm\varphi_{H'} = \norm a .
\]
\end{theorem}

\begin{proof}
यदि $\varphi = 0$ हो: $a = 0$। अन्यथा $F = \ker\varphi$ एक संवृत उचित उपसमष्टि है; $u \in F^\perp$,
$\norm u = 1$ चुनिए (\cref{thm:b3:hilbert:decomposition}: $F^\perp \neq 0$, क्योंकि $F \neq H$)। किसी भी $x$ के लिए सदिश
$\varphi(x)u -
\varphi(u)x \in \ker\varphi$, अतः $\perp u$:
\[
0 = \langle u, \varphi(x)u - \varphi(u)x\rangle
= \varphi(x) - \varphi(u)\langle u, x\rangle :
\qquad \varphi(x) = \langle
\overline{\varphi(u)}\,u,\ x\rangle .
\]
इसलिए $a = \overline{\varphi(u)}u$ काम कर जाता है। अद्वितीयता: सभी $x$ के लिए $\langle a
- a', x\rangle = 0$;
$x = a - a'$ की परीक्षा लीजिए। मानक: $\abs{\varphi(x)} \leq \norm a\norm x$ (कोशी--श्वार्ज़), और $x = a$ पर
समता।
\end{proof}

\begin{example}[एक प्रक्षेप, अंत तक परिकलित]
\label{ex:b3:hilbert:projectionexample}
$H = L^2(\intcc01)$ में $f(x) = x^2$ का किसी आफ़ीन फलन से श्रेष्ठतम आसन्नन क्या है?
उपसमष्टि $F =
\operatorname{Vect}(1, x)$ संवृत है (परिमित-विमीय), और $p_F(f) = a + bx$ इस गुण से
अभिलक्षित होता है कि अवशेष $1$ और $x$ दोनों के लांबिक है:
\[
\int_0^1(x^2 - a - bx)\,\dd x = 0,
\qquad
\int_0^1x\,(x^2 - a - bx)\,\dd x = 0,
\]
अर्थात् $\frac13 = a + \frac b2$ और $\frac14 = \frac a2 +
\frac b3$: अतः $a = -\frac16$, $b = 1$। इसलिए $p_F(x^2) = x -
\frac16$, और त्रुटि
\[
d(f, F)^2 = \int_0^1\Bigl(x^2 - x + \frac16\Bigr)^2\dd x =
\frac1{180},
\qquad d(f, F) = \frac1{6\sqrt5} .
\]
है। दो बातें आत्मसात करने योग्य हैं। पहली, यह परिकलन एक $2\times2$
रैखिक निकाय के अतिरिक्त कुछ नहीं है --- \emph{प्रसामान्य समीकरण};
एकपदी आधार में उनका आव्यूह $\bigl(\frac1{i+j+1}\bigr)$ कुख्यात रूप से दुर्दशित हिल्बर्ट
आव्यूह है, और पहले लांबिकीकरण कर लेना (लजांद्र बहुपद, \cref{pb:b3:hilbert:1}) उसका
उपचार है। दूसरी, $x^2$ का आफ़ीन फलनों से श्रेष्ठतम \emph{एकसमान}
आसन्नन भिन्न है ($x - \frac18$, समदोलन से): प्रत्येक मानक की अपनी ज्यामिति
होती है, और केवल हिल्बर्टीय मानक ही किसी रैखिक निकाय से उत्तर देता
है।
\end{example}

\section{प्रसामान्य लांबिक आधार}

\begin{definition}\label{def:b3:hilbert:onb}
कोई कुल $(e_i)_{i\in I}$ \emph{प्रसामान्य लांबिक} कहलाता है यदि $\langle
e_i, e_j\rangle = \delta_{ij}$ हो,
और \emph{हिल्बर्ट आधार}\index{हिल्बर्ट आधार} (प्रसामान्य लांबिक
आधार) यदि इसके अतिरिक्त उसके परिमित रैखिक संचय $H$ में सघन हों
(अर्थात् कुल \emph{संपूर्ण} हो)। हम गणनीय स्थिति $I = \N$ लेते हैं,
जो ग्राम--श्मिट के माध्यम से प्रत्येक \emph{पृथक्करणीय} $H$ को
समेट लेती है (\cref{prop:b3:hilbert:gramschmidt})।
\end{definition}

\begin{theorem}[बेसेल, पार्सेवाल]\label{thm:b3:hilbert:parseval}
मान लीजिए $(e_n)_{n\in\N}$ $H$ में प्रसामान्य लांबिक है और $c_n(x) =
\langle e_n, x\rangle$।
\begin{enumerate}
  \item (बेसेल) $\sum_n\abs{c_n(x)}^2 \leq \norm x^2$, और श्रेणी $\sum_nc_n(x)e_n$ $H$ में अभिसरित होती है,
        जिसका योग $p_F(x)$, $F = \overline{\operatorname{Vect}}(e_n)$ है।
  \item निम्नलिखित समतुल्य हैं: (क) $(e_n)$ एक \omterm{def:b3:hilbert:onb}{हिल्बर्ट आधार} है;
        (ख) प्रत्येक $x$ के लिए $x = \sum_nc_n(x)e_n$; (ग)
        \emph{पार्सेवाल}\index{पार्सेवाल की सर्वसमिका}: प्रत्येक
        $x$ के लिए $\norm x^2
        = \sum_n\abs{c_n(x)}^2$; (घ) समस्त $e_n$ के लांबिक एकमात्र
        सदिश $0$ है।
  \item यदि $(e_n)$ एक \omterm{def:b3:hilbert:onb}{हिल्बर्ट आधार} हो, तो $x \mapsto (c_n(x))_n$ एक समदूरिक
        तुल्याकारिता $H \to \ell^2$ है (\emph{प्रत्येक} अनंत-विमीय
        पृथक्करणीय \omterm{def:b3:hilbert:inner}{हिल्बर्ट समष्टि} ``$\ell^2$ ही है''), और $\langle x, y\rangle =
        \sum_n\overline{c_n(x)}c_n(y)$।
\end{enumerate}
\end{theorem}

\begin{proof}
(1) परिमित $N$ के लिए: $x - \sum_{n\leq N}c_ne_n \perp e_k$ ($k
\leq N$), अतः पाइथागोरस से $\norm x^2 = \sum_{n\leq
N}\abs{c_n}^2 + \norm{x - \sum_{n\leq N}c_ne_n}^2$:
यही बेसेल है। आंशिक योग $S_N = \sum_{n\leq N}c_ne_n$ कोशी हैं: $\norm{S_N - S_M}^2 = \sum_{M<n\leq N}\abs{c_n}^2$, जो किसी अभिसारी
श्रेणी की पुच्छ है; सीमा $F$ में होती है, और \omterm{def:b3:topology:continuity}{सांतत्य} से प्रत्येक
$e_k$ पर $x - \lim S_N
\perp$, अतः $\perp F$: लांबिक विघटन की अद्वितीयता से $\lim S_N = p_F(x)$।

(2) (क)$\Rightarrow$(ख): $F = H$, अतः $p_F = \mathrm{id}$।
(ख)$\Rightarrow$(ग): सीमा में पाइथागोरस ($\norm{S_N}^2
= \sum_{n \leq N}\abs{c_n}^2 \to \norm x^2$)।
(ग)$\Rightarrow$(घ): सभी $e_n$ के लिए $x \perp$ से $\norm x^2 =
0$ मिलता है।
(घ)$\Rightarrow$(क): $F^\perp = \{0\}$ (समस्त $e_n$ के लांबिक होना $F$ के लांबिक
होना है), अतः \cref{thm:b3:hilbert:decomposition} से $F$ सघन है; परंतु $F$, जो एक \omterm{def:b3:topology:interior}{संवृति}
है, पहले से ही संवृत है: अतः $F = H$।

(3) प्रतिचित्रण रैखिक है, (ग) से समदूरिक (अतः एकैकी), और आच्छादक
भी: $(c_n) \in \ell^2$ दिया हो तो श्रेणी $\sum
c_ne_n$ ((1) की तरह कोशी) किसी
पूर्वप्रतिबिंब पर अभिसरित होती है। \omterm{def:b3:hilbert:inner}{अंतःगुणन} सूत्र (ग) से ध्रुवण
द्वारा, अथवा किसी सीधे सीमा परिकलन से मिलता है।
\end{proof}

\begin{proposition}[ग्राम--श्मिट]\label{prop:b3:hilbert:gramschmidt}
मान लीजिए $(x_n)$ एक रैखिकतः स्वतंत्र अनुक्रम है। आगमन से $\tilde e_n = x_n - \sum_{k<n}\langle e_k,
x_n\rangle e_k$ और
$e_n = \tilde e_n/\norm{\tilde e_n}$ रखने पर एक प्रसामान्य लांबिक $(e_n)$ मिलता है जिसके परिमित
विस्तार वही हैं: $\operatorname{Vect}(e_1, \dots, e_n) = \operatorname{Vect}
(x_1, \dots, x_n)$। फलस्वरूप प्रत्येक पृथक्करणीय \omterm{def:b3:hilbert:inner}{हिल्बर्ट
समष्टि} (जिसका कोई गणनीय सघन उपसमुच्चय हो) का कोई \omterm{def:b3:hilbert:onb}{हिल्बर्ट आधार}
होता है।\index{ग्राम--श्मिट प्रसामान्य लांबिकीकरण}
\end{proposition}

\begin{proof}
आगमन: रचना से $\tilde e_n \perp e_k$ ($k < n$), और स्वतंत्रता से $\tilde e_n \neq 0$; प्रत्येक चरण
पर विस्तार मेल खाते हैं (त्रिभुजाकार आधार-परिवर्तन)। पृथक्करणीय
$H$ के लिए: किसी सघन अनुक्रम से सघन विस्तार वाला कोई रैखिकतः
स्वतंत्र उपकुल निकालिए (हर उस सदिश को छोड़ दीजिए जो अपने पूर्ववर्तियों
के विस्तार में हो --- विस्तार अपरिवर्तित रहता है), और उसका
प्रसामान्य लांबिकीकरण कीजिए: परिणाम संपूर्ण होता है।
\end{proof}

\begin{theorem}[त्रिकोणमितीय निकाय; अंततः पार्सेवाल]
\label{thm:b3:hilbert:fourier}
$\langle f, g\rangle =
\frac1{2\pi}\int_{-\pi}^\pi \bar fg$ के साथ $L^2(\intcc{-\pi}\pi)$ में कुल $e_n(t) =
\eu^{\iu nt}$, $n \in \Z$, एक \omterm{def:b3:hilbert:onb}{हिल्बर्ट आधार} है।
फलस्वरूप \emph{प्रत्येक} $f \in L^2$ के लिए --- विशेष रूप से प्रत्येक
खंडशः \omterm{def:b3:topology:continuity}{संतत} $2\pi$-आवर्ती $f$ के लिए --- $c_n(f) =
\frac1{2\pi}\int_{-\pi}^{\pi}f(t)\eu^{-\iu nt}\dd t$ के साथ:
\[
f = \sum_{n\in\Z}c_n(f)\,\eu^{\iu nt} \ \ L^2 \text{ में},
\qquad
\frac1{2\pi}\int_{-\pi}^{\pi}\abs f^2 =
\sum_{n\in\Z}\abs{c_n(f)}^2 .
\]
यह उसी \omterm{thm:b3:hilbert:parseval}{पार्सेवाल} सर्वसमिका को पूरी व्यापकता में सिद्ध कर देता है
जिसे दूसरे वर्ष ने स्वीकार कर लिया था।\index{पार्सेवाल की
सर्वसमिका (व्यापक)}
\end{theorem}

\begin{proof}
प्रसामान्य लांबिकता एक सीधा परिकलन है (दूसरा वर्ष)। संपूर्णता: मान
लीजिए $f \in L^2$ समस्त $e_n$ के $\perp$ है, अर्थात् उसके सभी फूरिये
गुणांक लुप्त हैं। \omterm{def:b3:topology:continuity}{संतत} $2\pi$-आवर्ती फलन $L^2(\intcc{-\pi}\pi)$ में सघन हैं: वस्तुतः
$\mathcal
C_c(\intoo{-\pi}\pi)$ सघन है (\cref{thm:b3:lp:density}(2)) और ऐसे फलन आवर्ती तथा \omterm{def:b3:topology:continuity}{संतत} रूप से
विस्तारित हो जाते हैं। त्रिकोणमितीय बहुपद \omterm{def:b3:topology:continuity}{संतत} आवर्ती फलनों में
$\norm\cdot_\infty$-सघन हैं (स्टोन--वाइरश्ट्रास, \cref{cor:b3:complete:weierstrass}(ग)), और $\norm\cdot_2 \leq \norm\cdot_\infty$: अतः
त्रिकोणमितीय बहुपद $L^2$ में सघन हैं। परंतु $f \perp$ प्रत्येक
त्रिकोणमितीय बहुपद, अतः $f \perp$ किसी सघन उपसमष्टि: इसलिए $f
\in (\text{सघन})^\perp = \{0\}$
(\cref{thm:b3:hilbert:decomposition})। \cref{thm:b3:hilbert:parseval} की कसौटी (घ) निष्कर्ष दे देती है; और (ख) तथा (ग)
खुलकर प्रदर्शित रूप बन जाते हैं (गणनीय $\Z$ का पुनःसूचकांकन; और
दो-मुखी श्रेणी अप्रतिबंधित रूप से अभिसरित होती है --- किसी भी
निःशेषक कुल पर आंशिक योग $\ell^2$ पुच्छ तर्क से अभिसरित होते हैं)।
\end{proof}

\begin{theorem}[लाक्स--मिलग्राम]\label{thm:b3:hilbert:laxmilgram}
मान लीजिए $H$ एक वास्तविक \omterm{def:b3:hilbert:inner}{हिल्बर्ट समष्टि} है और $a \colon H\times H \to
\R$
द्विरैखिक, \emph{\omterm{def:b3:topology:continuity}{संतत}} ($\abs{a(u,v)} \leq M\norm
u\norm v$) तथा \emph{प्रबलकारी} ($a(u, u) \geq \alpha\norm u^2$,
$\alpha > 0$)। तब प्रत्येक $\varphi \in H'$ के लिए एक अद्वितीय $u \in H$ है जिसके
लिए\index{लाक्स--मिलग्राम प्रमेय}
\[
a(u, v) = \varphi(v) \qquad \text{प्रत्येक } v \in H \text{ के लिए} .
\]
\end{theorem}

\begin{proof}
स्थिर $u$ के लिए $v \mapsto a(u, v)$ \omterm{def:b3:topology:continuity}{संतत} रैखिक है: रीस एक अद्वितीय $Au \in H$ दे
देता है जिसके लिए $a(u,v) = \langle Au,
v\rangle$; और $A$ $\norm{Au} \leq M\norm u$ के साथ रैखिक है
(प्रतिनिधियों की अद्वितीयता, फिर परिबंध)। प्रबलकारिता: $\alpha\norm u^2 \leq a(u,u) = \langle Au, u\rangle \leq
\norm{Au}\norm u$, अतः
$\norm{Au} \geq \alpha\norm u$: इसलिए $A$ एकैकी है और उसका परिसर संवृत है (कोई कोशी
प्रतिबिंब अनुक्रम $Au_n$ $u_n$ को कोशी होने पर विवश कर देता है)।
परिसर सघन है: $w \perp
\operatorname{im}A$ से $0 = \langle Aw, w\rangle \geq
\alpha\norm w^2$ मिलता है। संवृत और सघन: अतः $A$
एकैकी आच्छादक है। $\varphi$ दिया हो, तो उसे निरूपित करने वाला $f$
लीजिए (रीस) और $u = A^{-1}f$: तब $a(u, v) = \langle f, v\rangle = \varphi(v)$, और वह भी अद्वितीय रूप से ($a(u - u', \cdot) = 0$
और प्रबलकारिता)।
\end{proof}

\begin{remark}\label{rem:b3:hilbert:variational}
जब $a$ सममित हो, तब लाक्स--मिलग्राम का हल \emph{ऊर्जा} $J(v) = \frac12a(v,v) -
\varphi(v)$
(\cref{exo:b3:hilbert:9}) का अद्वितीय न्यूनीकारक होता है: अर्थात् एक ही आघात में
विचरण समस्याओं के हलों का अस्तित्व। उपयुक्त फलन समष्टियों (आगे के
पाठ्यक्रम की सोबोलेव समष्टियाँ) पर लगाकर यह अवकल समीकरणों की सीमांत
मान समस्याएँ हल कर देता है --- आंशिक अवकल समीकरणों में आधुनिक
प्रवेश द्वार।
\end{remark}

\section{अभ्यास}

\begin{exercise}[$\star$]\label{exo:b3:hilbert:1}
(क) ध्रुवण सर्वसमिकाएँ सिद्ध कीजिए (वास्तविक: $4\langle x,
y\rangle = \norm{x+y}^2 - \norm{x-y}^2$; सम्मिश्र:
चार-पद वाला रूप)।
(ख) दर्शाइए कि $L^1(\intcc01)$ पर $\norm\cdot_1$ और $\mathcal C(\intcc01)$ पर $\norm\cdot_\infty$ समांतर चतुर्भुज
नियम का उल्लंघन करते हैं: अतः ये मानक किसी \omterm{def:b3:hilbert:inner}{अंतःगुणन} से नहीं आते।
\end{exercise}

\begin{exercise}[$\star$]\label{exo:b3:hilbert:2}
$H = L^2(\intcc01)$ (वास्तविक) में:
(क) अचर फलनों की उपसमष्टि पर $f$ का प्रक्षेप परिकलित कीजिए और
उसकी व्याख्या कीजिए;
(ख) $\{g : g = 0 \text{, }
\intcc0{1/2} \text{ पर लगभग सर्वत्र}\}$ पर प्रक्षेप परिकलित कीजिए;
(ग) $d\bigl(x \mapsto x,\ \operatorname{Vect}(\mathbf
1)\bigr)$ परिकलित कीजिए।
\end{exercise}

\begin{exercise}[$\star\star$]\label{exo:b3:hilbert:3}
(क) दर्शाइए कि किसी उपसमष्टि $F$ के लिए: $F$ सघन $\iff$ $F^\perp =
\{0\}$,
और $\ell^2$ में किसी \emph{उचित} सघन उपसमष्टि का उदाहरण दीजिए (जिससे
$F = H$ के बिना $F^\perp = 0$: विघटन प्रमेय को $F$ का संवृत होना वास्तव
में चाहिए)।
(ख) दर्शाइए कि यदि $x_n \to x$ और मानक में $y_n \to y$ हो, तो $\langle x_n, y_n\rangle \to \langle x, y\rangle$, और इस
अध्याय में दो ऐसे स्थान बताइए जहाँ इस \omterm{def:b3:topology:continuity}{सांतत्य} का उपयोग हुआ।
\end{exercise}

\begin{exercise}[$\star\star$]\label{exo:b3:hilbert:4}
$L^2(\intcc{-1}1)$ (\omterm{def:b3:measure:lebesgueouter}{लेबेग माप}) में $1, x, x^2$ पर ग्राम--श्मिट लगाइए: पहले तीन
मानकीकृत \emph{लजांद्र बहुपद} प्राप्त कीजिए, और सत्यापित कीजिए कि
वे \cref{pb:b3:hilbert:1} के रोद्रीग बहुपदों $P_n$ के लिए $\sqrt{n + \frac12}\,P_n$ से मेल खाते हैं।
\end{exercise}

\begin{exercise}[$\star\star$]\label{exo:b3:hilbert:5}
$\intcc{-\pi}\pi$ पर $f(t) = t$ और $f(t) = t^2$ पर \omterm{thm:b3:hilbert:parseval}{पार्सेवाल} (\cref{thm:b3:hilbert:fourier}) लगाइए --- अब इनके लिए
वैध रूप से (ये \omterm{def:b3:topology:continuity}{संतत} हैं, पर पहले सर्वसमिका को लपेटन-बिंदु की असांतत्य
पर खंडशः-$\mathcal C^1$ सावधानी चाहिए थी): और
\[
\sum_{n\geq1}\frac1{n^2} = \frac{\pi^2}6,
\qquad
\sum_{n\geq1}\frac1{n^4} = \frac{\pi^4}{90} .
\]
पुनः प्राप्त कीजिए।
\end{exercise}

\begin{exercise}[$\star\star$]\label{exo:b3:hilbert:6}
(क) ऐसा $a \in L^2(\intcc01)$ ढूँढ़िए जिसके लिए सभी $f$ पर $\int_0^{1/2}f =
\langle a, f\rangle$ हो; इस फलनिक
के लिए $\norm\varphi$ परिकलित कीजिए।
(ख) दर्शाइए कि उपसमष्टि $\mathcal C(\intcc01) \subseteq
L^2(\intcc01)$ पर परिभाषित मान-निर्धारण $f \mapsto f(\frac12)$
$\norm\cdot_2$ के लिए \omterm{def:b3:topology:continuity}{संतत} \emph{नहीं} है: अतः कोई रीस प्रतिनिधि विद्यमान
नहीं (मान-निर्धारण कोई $L^2$ धारणा नहीं है)।
\end{exercise}

\begin{exercise}[$\star\star\star$]\label{exo:b3:hilbert:7}
मान लीजिए $H$ \omterm{def:b3:hilbert:onb}{हिल्बर्ट आधार} $(e_n)$ के साथ पृथक्करणीय है और $(x_k)$
कोई परिबद्ध अनुक्रम।
(क) दर्शाइए कि कोई उपानुक्रम \emph{दुर्बल} रूप से अभिसरित होता है:
ऐसा $x$ है कि प्रत्येक $y \in H$ के लिए $\langle y, x_{k_j}\rangle \to \langle y,
x\rangle$। \emph{(गुणांकों
$\langle e_n, x_k\rangle$ पर विकर्ण निष्कर्षण; बेसेल और मानकों की एकसमान परिबद्धता के
माध्यम से $x$ संकलित कीजिए।)}
(ख) दर्शाइए कि $e_n \rightharpoonup 0$ पर $\norm{e_n} = 1$: दुर्बल सीमाएँ मानक खो सकती हैं।
(क) में $\norm x \leq
\liminf\norm{x_{k_j}}$ दर्शाइए।
\end{exercise}

\begin{exercise}[$\star\star$]\label{exo:b3:hilbert:8}
(संलग्न) $T \in \mathcal L(H)$ के लिए दर्शाइए कि $\langle Tx, y\rangle = \langle
x, T^*y\rangle$ वाला एक अद्वितीय $T^* \in \mathcal L(H)$ है
(रीस), और $\vertiii{T^*} = \vertiii T$। $\ell^2$ पर विस्थापन $S$ का संलग्न परिकलित कीजिए,
और $\ker T^* = (\operatorname{im}T)^\perp$ सिद्ध कीजिए --- इससे $\overline{\operatorname{im}T} = (\ker T^*)^\perp$ निष्कर्ष निकालिए।
\end{exercise}

\begin{exercise}[$\star\star$]\label{exo:b3:hilbert:9}
मान लीजिए $a$ लाक्स--मिलग्राम जैसा है और इसके अतिरिक्त
\emph{सममित} भी। दर्शाइए कि $u$ $a(u, \cdot) = \varphi$ को हल करता है तभी और केवल
तभी जब $u$ $J(v) = \frac12a(v, v) - \varphi(v)$ को न्यूनतम करे, और यह कि न्यूनतम ठीक एक ही
बिंदु पर प्राप्त होता है। \emph{(वर्ग पूरा कीजिए: $J(u + w) - J(u) = \frac12a(w,w) \geq
\frac\alpha2\norm w^2$।)}
अनुप्रयोग: लाक्स--मिलग्राम से संवृत उपसमष्टियों के लिए प्रक्षेप
प्रमेय पुनः व्युत्पन्न कीजिए।
\end{exercise}

\begin{exercise}[$\star\star\star$]\label{exo:b3:hilbert:10}
(हार निकाय) $\intcc01$ पर मान लीजिए $h_{0} = \mathbf 1$, और $n = 2^j + k$ ($j \geq 0$, $0 \leq k < 2^j$) के
लिए:
\[
h_n = 2^{j/2}\Bigl(\mathbf 1_{[k2^{-j},\,(k +
\frac12)2^{-j})} - \mathbf 1_{[(k+\frac12)2^{-j},\,(k+1)2^{-j})}
\Bigr).
\]
दर्शाइए कि $(h_n)_{n\geq0}$ $L^2(\intcc01)$ में प्रसामान्य लांबिक और संपूर्ण है।
\emph{(लांबिकता: असंयुक्त या नेस्टेड आलंब; संपूर्णता: परिमित विस्तार
समस्त द्विआधारी सोपान फलनों को समेटते हैं, जो सघन हैं --- \cref{thm:b3:lp:density}(1)
और अंतरालों के द्विआधारी आसन्नन के माध्यम से।)} हार निकाय
तरंगिकाओं का पूर्वज है।
\end{exercise}

\begin{exercise}[$\star\star$]\label{exo:b3:hilbert:11}
(लांबिक प्रक्षेप, अभिलक्षित) मान लीजिए $H$ एक \omterm{def:b3:hilbert:inner}{हिल्बर्ट समष्टि} है
और $P^2 = P$, $P \neq 0$ वाला $P \in \mathcal L(H)$। इनकी तुल्यता दर्शाइए:
(क) $P$ $\operatorname{im}P$ पर लांबिक प्रक्षेप है;
(ख) $P = P^*$ (\cref{exo:b3:hilbert:8});
(ग) $\vertiii P = 1$।
\emph{((ग) $\Rightarrow$ (क) के लिए: यदि किसी $x \in
(\ker P)^\perp$ के लिए $Px \neq x$ हो, तो
$x + t(Px - x)$ पर विचार कीजिए --- अथवा सीधे: $u \in \operatorname{im}P$ और $v \in
\ker P$ के लिए सभी
$t \in \R$ पर $\norm{P(u + tv)}^2 \leq \norm{u + tv}^2$ का प्रसार कीजिए और $\langle u, v\rangle = 0$ निष्कर्ष निकालिए।)}
$\R^2$ पर कोई अलांबिक प्रक्षेप प्रस्तुत कीजिए और उसका मानक परिकलित
कीजिए।
\end{exercise}

\begin{exercise}[$\star\star\star$]\label{exo:b3:hilbert:12}
(फोन नॉयमान की एर्गोडिक प्रमेय) मान लीजिए $U \in \mathcal L(H)$ \emph{एकात्मक}
है ($U^*U = UU^* = I$), $F = \ker(U - I)$ स्थिर समष्टि है, $P$ $F$ पर लांबिक प्रक्षेप,
और $A_n = \frac1n\sum_{k=0}^{n-1}U^k$।
(क) दर्शाइए कि $\ker(U - I) = \ker(U^* - I)$ \emph{($\norm{Ux - x}^2 = 2\norm x^2 - 2\operatorname{Re}\langle
Ux, x\rangle$ और एकात्मकता से)}, और $\overline{\operatorname{im}(U - I)} = F^\perp$
निष्कर्ष निकालिए।
(ख) दर्शाइए कि $x \in F$ के लिए $A_nx \to x$, और $x \in \operatorname{im}(U - I)$ के लिए $A_nx \to 0$
\emph{(दूरबीनी योग)}, फिर $x \in \overline{\operatorname{im}(U - I)}$ के लिए भी (एकसमान परिबंध $\vertiii{A_n} \leq 1$)।
(ग) निष्कर्ष निकालिए: \emph{प्रत्येक} $x \in H$ के लिए $A_nx \to Px$ ---
अर्थात् समय-औसत अपरिवर्तियों पर के प्रक्षेप पर अभिसरित होते हैं।
(घ) $H = L^2(\R/\Z)$ और अपरिमेय $\alpha$ वाले $Uf = f(\cdot +
\alpha)$ के लिए इसे विस्तार से लिखिए:
$F$ पहचानिए (फूरिये श्रेणी का उपयोग कीजिए, \cref{thm:b3:hilbert:fourier}) और निष्कर्ष
निकालिए कि $L^2$ में $\frac1n\sum_{k<n}f(x + k\alpha) \to \int_0^1f$: अर्थात् अपरिमेय घूर्णनों का $L^2$
समवितरण।
\end{exercise}

\section{समस्या: लांबिक बहुपद}

\begin{problem}[{सप्ताहांत समस्या --- लजांद्र, एर्मीत, और गाउस
संख्यात्मक समाकलन}]\label{pb:b3:hilbert:1}
मान लीजिए $I \subseteq \R$ कोई अंतराल है और $w > 0$ $I$ के \omterm{def:b3:topology:interior}{अंतःभाग} पर ऐसा
\omterm{def:b3:topology:continuity}{संतत} \emph{भार} कि सभी $n$ के लिए $\int_I
\abs t^nw(t)\dd t < \infty$; $\langle f, g\rangle = \int_I \bar
fg\,w$ के साथ $H = L^2(I,
w\,\dd\lambda)$ में
कार्य कीजिए। $1, t, t^2, \dots$ पर ग्राम--श्मिट लगाने से $w$ के लिए
\emph{\omterm{pb:b3:hilbert:1}{लांबिक बहुपद}} $(p_n)$ मिलते हैं (मोनिक मानकीकरण:
$p_n = t^n + \cdots$)।\index{लांबिक बहुपद}

\medskip
\noindent\textbf{भाग I --- व्यापक सिद्धांत।}
\begin{enumerate}
  \item दर्शाइए कि $p_n$ घात $< n$ के प्रत्येक बहुपद के लांबिक
        है, और यह कि $(p_0, \dots, p_n)$ $\R_n[t]$ का आधार है।
  \item (त्रिपद पुनरावर्तन) दर्शाइए कि ऐसे वास्तविक $a_n,
        b_n$ हैं
        जिनके लिए
        \[
        p_{n+1}(t) = (t - a_n)\,p_n(t) - b_n\,p_{n-1}(t),
        \qquad b_n = \frac{\norm{p_n}^2}{\norm{p_{n-1}}^2} >
        0 .
        \]
        \emph{($t\,p_n$ का आधार $(p_k)_{k \leq
        n+1}$ में प्रसार कीजिए और $\langle tp_n, p_k\rangle = \langle p_n,
        tp_k\rangle$ का
        उपयोग करते हुए लांबिकता से गुणांक समाप्त कीजिए।)}
  \item (मूल) दर्शाइए कि $p_n$ के $n$ \emph{भिन्न} मूल हैं, और
        सभी $I$ के \omterm{def:b3:topology:interior}{अंतःभाग} में। \emph{(मान लीजिए $t_1 < \dots < t_m$ $p_n$ के
        आंतरिक चिह्न-परिवर्तन हैं; यदि $m < n$ हो, तो $\prod_{i\leq m}(t - t_i)$ के
        सापेक्ष $p_n$ की परीक्षा लीजिए और लांबिकता का खंडन
        कीजिए।)}
\end{enumerate}

\medskip
\noindent\textbf{भाग II --- लजांद्र ($I = \intcc{-1}1$, $w =
1$)।} $P_n(t) = \frac{1}{2^nn!}\,\frac{\dd^n}{\dd
t^n}\bigl[(t^2 - 1)^n\bigr]$
(रोद्रीग) परिभाषित कीजिए।
\begin{enumerate}[resume]
  \item दर्शाइए कि $\deg P_n = n$ का अग्र गुणांक $\frac{(2n)!}{2^n(n!)^2}$ है, और $n$ बार
        खंडशः समाकलन करके यह भी कि घात $< n$ के प्रत्येक बहुपद
        $Q$ के लिए $\langle P_n, Q\rangle = 0$: अर्थात् $P_n$ (मानकीकरण तक) $w = 1$ के
        लिए \omterm{pb:b3:hilbert:1}{लांबिक बहुपद} हैं।
  \item $\norm{P_n}_2^2 = \frac{2}{2n+1}$ परिकलित कीजिए \emph{($n$ बार अपने ही सापेक्ष
        खंडशः समाकलन कीजिए और किसी बीटा/वालिस समाकल तक सिमटाइए,
        \cref{exo:b3:product:8})}।
  \item दर्शाइए कि मानकीकृत लजांद्र बहुपद $L^2(\intcc{-1}1)$ का \omterm{def:b3:hilbert:onb}{हिल्बर्ट आधार}
        बनाते हैं \emph{(वाइरश्ट्रास, \cref{cor:b3:complete:weierstrass}, और $L^2$ में $\mathcal C$
        की सघनता)}, और $f(t) = \abs t$ का घात $2$ तक प्रसार कीजिए: $\abs t$
        का श्रेष्ठतम द्विघात $L^2$-आसन्नन परिकलित कीजिए।
\end{enumerate}

\medskip
\noindent\textbf{भाग III --- एर्मीत ($I = \R$, $w(t) =
\eu^{-t^2}$)।} $H_n(t) =
(-1)^n\eu^{t^2}\frac{\dd^n}{\dd t^n}\eu^{-t^2}$
परिभाषित कीजिए।
\begin{enumerate}[resume]
  \item दर्शाइए कि $H_n$ घात $n$ का ऐसा बहुपद है जिसका अग्र
        गुणांक $2^n$ है, कि $H_{n+1} = 2tH_n -
        H_n'$, और यह कि $\langle H_m, H_n\rangle_w =
        \delta_{mn}\,2^nn!\sqrt\pi$ \emph{(फिर
        खंडशः)}।
  \item दर्शाइए कि एर्मीत कुल $L^2(\R,
        \eu^{-t^2}\dd t)$ में संपूर्ण है, और इसके लिए
        \cref{ch:b3:fouriertransform} का एक परिणाम स्वीकार कीजिए: यदि $g \in L^1(\R)$ के लिए सभी
        $\xi$ पर $\int g(t)\eu^{-\iu\xi t}\dd t = 0$ हो, तो लगभग सर्वत्र $g = 0$। \emph{(सभी
        $H_n$ के लिए $f \perp$, अर्थात् समस्त बहुपदों के $\perp$: दर्शाइए
        कि $z \mapsto \int
        f(t)\eu^{-t^2}\eu^{-\iu zt}\dd t$ सुपरिभाषित है, चरघातांकी का श्रेणी में प्रसार
        कीजिए, प्रभुत्व से अदला-बदली उचित ठहराइए, और निष्कर्ष
        निकालिए कि $f\eu^{-t^2}$ का फूरिये रूपांतर लुप्त हो जाता है।)}
\end{enumerate}

\medskip
\noindent\textbf{भाग IV --- गाउस संख्यात्मक समाकलन।} $n$ स्थिर
कीजिए, मान लीजिए $t_1 < \dots < t_n$ $p_n$ (भाग I) के मूल हैं, और भार $w_i = \int_I \ell_i(t)\,w(t)\dd t$
परिभाषित कीजिए, जहाँ $\ell_i$ $t_i$ पर लाग्राँज अंतर्वेशन आधार बहुपद
हैं।
\begin{enumerate}[resume]
  \item दर्शाइए कि संख्यात्मक समाकलन नियम $Q(f) = \sum_iw_if(t_i)$ घात $\leq n - 1$ तक के
        सभी बहुपदों पर यथार्थ है (अंतर्वेशन), और वस्तुतः ---
        यही चमत्कार है --- घात $\leq 2n - 1$ तक के सभी बहुपदों पर: $P =
        qp_n + r$
        लिखिए और भागफल $q$ पर लांबिकता का उपयोग कीजिए।
  \item दर्शाइए कि भार धनात्मक हैं \emph{(नियम को घात $2n -
        2$ के
        $\ell_i^2$ पर लगाइए)}, और पोल्या की प्रमेय (\cref{exo:b3:banach:9}) से निष्कर्ष
        निकालिए कि गाउस संख्यात्मक समाकलन अभिसरित होता है: किसी
        \omterm{def:b3:topology:compact}{संहत} $I$ पर प्रत्येक \omterm{def:b3:topology:continuity}{संतत} $f$ के लिए $Q_n(f) \to \int_I fw$।
  \item $n = 2$, $I = \intcc{-1}1$, $w = 1$ के लिए: नोड $\pm\frac1{\sqrt3}$ और भार $1, 1$
        परिकलित कीजिए, और $1, t, t^2, t^3$ पर यथार्थता हाथ से सत्यापित
        कीजिए। उन्हीं दो मान-निर्धारण बिंदुओं पर समलंब नियम से
        तुलना कीजिए।
\end{enumerate}

\medskip
\noindent\textbf{भाग V --- चेबिशेव: सबसे अच्छा दोलन करने वाले
बहुपद।} अब $I = \intcc{-1}1$ और $w(t) =
\frac1{\sqrt{1 - t^2}}$।
\begin{enumerate}[resume]
  \item दर्शाइए कि $T_n(\cos\theta) = \cos n\theta$ घात $n$ का एक बहुपद $T_n$ परिभाषित
        करता है (किसी त्रिकोणमितीय सर्वसमिका से $T_{n+1} =
        2t\,T_n - T_{n-1}$ स्थापित
        कीजिए), जिसका अग्र गुणांक $n \geq 1$ के लिए $2^{n-1}$ है; और यह कि
        प्रतिस्थापन $t = \cos\theta$ से
        \[
        \langle T_m, T_n\rangle_w =
        \int_0^\pi\cos m\theta\,\cos n\theta\,\dd\theta
        = 0 \ (m \neq n), \qquad
        \norm{T_0}_w^2 = \pi,\ \ \norm{T_n}_w^2 = \frac\pi2 :
        \]
        मिलता है: अर्थात् $T_n$ इस भार के लिए \omterm{pb:b3:hilbert:1}{लांबिक बहुपद} हैं,
        और चेबिशेव प्रसार \emph{वस्तुतः} छद्म वेश में फूरिये
        कोज्या श्रेणियाँ हैं।
  \item $\intcc{-1}1$ पर $T_n$ के $n$ मूल $t_k =
        \cos\frac{(2k-1)\pi}{2n}$ और $n + 1$ चरम मान $s_j = \cos\frac{j\pi}n$
        स्पष्ट रूप से बताइए, जहाँ $T_n(s_j) = (-1)^j$: अर्थात् आलेख $\pm1$ के बीच
        \emph{समदोलन} करता है।
  \item (न्यूनतम-उच्चतम) दर्शाइए कि घात $n$ के समस्त
        \emph{मोनिक} बहुपदों में बहुपद $2^{1-n}T_n$ का $\intcc{-1}1$ पर उच्चतम
        मानक सबसे छोटा है, अर्थात् $2^{1-n}$ --- और वही अद्वितीय
        न्यूनीकारक है। \emph{(यदि किसी मोनिक $P$ के लिए $\sup\abs P < 2^{1-n}$
        होता, तो घात $\leq n-1$ का अंतर $2^{1-n}T_n -
        P$ $n+1$ समदोलन बिंदुओं पर
        चिह्न बदलता।)}
  \item अंतर्वेशन में अनुप्रयोग: $\intcc{-1}1$ में नोड $t_1 < \dots
        < t_n$ के लिए किसी
        $\mathcal C^n$ फलन के लाग्राँज अंतर्वेशन की त्रुटि में $\omega(t) = \prod_i(t - t_i)$ आता
        है। दर्शाइए कि नोड के रूप में चेबिशेव मूल चुनने से $\sup_{\intcc{-1}1}\abs\omega$
        न्यूनतम हो जाता है, और परिणामी परिबंध $\norm{f -
        L_nf}_\infty \leq \frac{\norm{f^{(n)}}_\infty}
        {2^{n-1}\,n!}$ दीजिए ---
        समान दूरी वाले नोड से तुलना कीजिए (चेतावनी की कथा के रूप
        में रुंगे परिघटना का कथन दीजिए)।
  \item $\abs{T_n'(\pm1)} = n^2$ सत्यापित कीजिए \emph{($T_n(\cos\theta) = \cos n\theta$ का अवकलन कीजिए और
        $\theta \to 0, \pi$ पर सीमाएँ लीजिए)}: अर्थात् $\intcc{-1}1$ पर $1$ से परिबद्ध
        बहुपदों का अवकलज किनारे पर $n^2$ जितना बड़ा हो सकता है
        (मार्कोव असमिका कहती है कि इससे बड़ा नहीं --- केवल कथन)।
        अंतराल में कहाँ अवकलज का परिबंध केवल $O(n)$ रहता है?
  \item (चेबिशेव--गाउस संख्यात्मक समाकलन) दर्शाइए कि भार $w$ के
        लिए $n$ चेबिशेव मूलों पर गाउस नियम के भार \emph{समान}
        होते हैं, $w_i = \frac\pi n$ \emph{($T_0, \dots, T_{n-1}$ पर यथार्थता और $1 \leq j \leq n - 1$ के लिए
        त्रिकोणमितीय योग $\sum_{k=1}^n\cos\bigl(j\tfrac{(2k-1)\pi}{2n}\bigr) =
        0$)}: अर्थात् सभी संख्यात्मक समाकलनों
        में सबसे एकसमान। $n = 3$ के लिए उसे लिखिए।
\end{enumerate}

\medskip
\noindent\textbf{भाग VI --- क्रिस्टोफेल--दारबू, अंतर्ग्रथन, और
याकोबी आव्यूह।} अब फिर किसी व्यापक भार पर लौटिए; $h_k = \norm{p_k}^2$ (मोनिक
$p_k$), $b_k = h_k/h_{k-1}$।
\begin{enumerate}[resume]
  \item (न्यूनतम मानक) दर्शाइए कि घात $n$ के समस्त \emph{मोनिक}
        बहुपदों में लांबिक $p_n$ ही न्यूनतम $L^2(w)$-मानक वाला
        अद्वितीय बहुपद है --- इस न्यूनीकरण को $\R_{n-1}[t]$ पर लांबिक
        प्रक्षेप के रूप में पहचानिए (\cref{thm:b3:hilbert:projection} अथवा दूसरे वर्ष का
        परिमित-विमीय प्रक्षेप)। प्रश्न 14 का न्यूनतम-उच्चतम गुण
        वही कथन है जिसमें $L^2$ के स्थान पर $L^\infty$ है: एक ही नायक,
        दो मानक।
  \item (क्रिस्टोफेल--दारबू) त्रिपद पुनरावर्तन का उपयोग करते हुए
        $n$ पर आगमन से सर्वसमिका
        \[
        \sum_{k=0}^{n}\frac{p_k(x)\,p_k(y)}{h_k}
        = \frac{p_{n+1}(x)\,p_n(y) -
        p_n(x)\,p_{n+1}(y)}{h_n\,(x - y)}
        \qquad (x \neq y),
        \]
        और उसका संगामी रूप ($y \to x$) सिद्ध कीजिए:
        $\sum_{k\leq n}\frac{p_k(x)^2}{h_k} =
        \frac{p_{n+1}'(x)p_n(x) - p_n'(x)p_{n+1}(x)}{h_n}$।
  \item इससे निष्कर्ष निकालिए कि $p_n$ और $p_{n+1}$ का कोई उभयनिष्ठ
        मूल नहीं है, और यह कि $p_{n+1}$ के प्रत्येक मूल $x_0$ पर
        $p_n(x_0)\,p_{n+1}'(x_0) > 0$। मूलों का \emph{अंतर्ग्रथन} निष्कर्ष के रूप में
        निकालिए: $p_{n+1}$ के दो क्रमागत मूलों के बीच $p_n$ का ठीक एक
        मूल होता है।
  \item (याकोबी आव्यूह) मान लीजिए $J_n$ वह $n\times n$ सममित त्रिविकर्ण
        आव्यूह है जिसका विकर्ण $a_0,
        \dots, a_{n-1}$ है और विकर्णेतर प्रविष्टियाँ
        $\sqrt{b_1},
        \dots, \sqrt{b_{n-1}}$। आगमन से दर्शाइए कि $\det(tI_n - J_n) = p_n(t)$, अतः $p_n$ के मूल किसी
        वास्तविक सममित आव्यूह के अभिलक्षणिक मान हैं --- जिससे एक
        पंक्ति में पुनः सिद्ध हो जाता है कि वे वास्तविक हैं, और
        (ऊपर के अंतर्ग्रथन के साथ) \omterm{pb:b3:hilbert:1}{लांबिक बहुपद} \cref{ch:b3:spectral} के
        स्पेक्ट्रमी संसार से जुड़ जाते हैं।
  \item (संश्लेषण) तीनों शास्त्रीय कुलों (लजांद्र, एर्मीत,
        चेबिशेव) के लिए शब्दकोश संकलित कीजिए: अंतराल, भार,
        परिभाषक सूत्र, त्रिपद पुनरावर्तन, मानक, और प्रत्येक का
        स्वाभाविक निवास (संहतों पर संख्यात्मक समाकलन और आसन्नन;
        गाउसीय विश्लेषण; न्यूनतम-उच्चतम और फूरिये-कोज्या विधियाँ)।
        एक वाक्य इस पर कि व्यापक सिद्धांत (भाग I, VI) ने वह क्या
        दिया जो कोई भी अलग परिकलन नहीं दे सकता था।
\end{enumerate}

\medskip
\noindent\textbf{भाग VII --- त्रुटि पद, और भारों के पीछे का
नाभिक।} यहाँ $I$ \omterm{def:b3:topology:compact}{संहत} है और $f \in \mathcal
C^{2n}(I)$।
\begin{enumerate}[resume]
  \item (गाउस त्रुटि सूत्र) मान लीजिए $Hf$ घात $\leq 2n - 1$ का वह
        एर्मीत अंतर्वेशक है जो नोड $t_1, \dots, t_n$ पर $f$ और $f'$ से मेल
        खाता है (उसका अस्तित्व और बिंदुशः त्रुटि
        \[
        f(t) - Hf(t) =
        \frac{f^{(2n)}(\xi_t)}{(2n)!}\;p_n(t)^2
        \]
        सामान्य सहायक-फलन तर्क से सिद्ध कीजिए)। इस सर्वसमिका का
        $w$ के सापेक्ष समाकलन करके और $f^{(2n)}$ के चरम मानों के
        बीच दबाकर निष्कर्ष निकालिए
        \[
        \int_I f\,w - Q_n(f)
        = \frac{f^{(2n)}(\xi)}{(2n)!}\,h_n
        \qquad\text{किसी } \xi \in I \text{ के लिए},
        \]
        जिसमें $h_n = \norm{p_n}^2$ भाग VI जैसा है: अर्थात् गाउस संख्यात्मक
        समाकलन एक $2n$-वें अवकलज जितनी त्रुटि करता है, और वह भी
        मोनिक \omterm{pb:b3:hilbert:1}{लांबिक बहुपद} के वर्ग-मानक से भारित।
  \item (भार क्रिस्टोफेल मान हैं) $\R_{n-1}[t]$ के पुनरुत्पादक नाभिक
        $K_n(x, y) = \sum_{k=0}^{n-1}
        \frac{p_k(x)p_k(y)}{h_k}$ और $Q_n$ की घात $2n - 2$ तक की यथार्थता का उपयोग करके
        सिद्ध कीजिए
        \[
        w_i \;=\;
        \Bigl(\,\sum_{k=0}^{n-1}
        \frac{p_k(t_i)^2}{h_k}\Bigr)^{\!-1} :
        \]
        अर्थात् प्रत्येक भार अपने ही नोड पर \emph{क्रिस्टोफेल
        फलन} का मान है --- भारों की धनात्मकता (प्रश्न 10) पुनः,
        अब एक यथार्थ सूत्र के साथ। सत्यापित कीजिए कि इससे $n =
        2$,
        $I = \intcc{-1}1$, $w = 1$ के लिए $w_1 = w_2 = 1$ पुनः मिल जाता है।
  \item (एक ही समाकल पर सब कुछ जाँचिए) चेबिशेव भार और $n = 3$ नोड
        के लिए
        \[
        \int_{-1}^{1}\frac{t^6}{\sqrt{1 - t^2}}\,\dd t
        = \frac{5\pi}{16},
        \qquad
        Q_3(t^6) = \frac{9\pi}{32},
        \]
        के दोनों पक्ष परिकलित कीजिए, जिससे संख्यात्मक समाकलन की
        त्रुटि ठीक $\frac{\pi}{32}$ निकले; फिर जाँचिए कि प्रश्न 23 का त्रुटि
        सूत्र ठीक यही मान बताता है (यहाँ $f^{(6)} = 6!$ अचर है, और $h_3 = \norm{2^{-2}T_3}_w^2 =
        \frac\pi{32}$):
        सिद्धांत और परिकलन अंतिम अंक तक सहमत हैं।
\end{enumerate}
\end{problem}
